\chapter{Methodology}
\label{ch:methodology}

\section{Overview}
\label{sec:methodology-overview}

This thesis formulates key--value pair extraction as two related
discriminative tasks. Entity classification assigns each observed text token
to \texttt{KEY}, \texttt{VALUE}, or \texttt{O}, where \texttt{O} denotes a
token that belongs to neither entity class. Relation linking then predicts
which \texttt{KEY} and \texttt{VALUE} candidates form a Regular key--value
pair.

The pipeline uses the pretrained LayoutLMv3-base encoder~\cite{layoutlmv3} on
KVP10k~\cite{kvp10k}. The pretrained encoder and the general
scaled-dot-product principle are prior work~\cite{layoutlmv3,vaswani2017attention}.
The task-specific classifier and linker modules, span-grouping path, training
pipeline, integration with the released evaluator, and diagnostic procedures
were implemented for this thesis. The released evaluator itself is a prior
resource.

The development variants differ in their relation candidates and training
paths. V1 is the token-level linker variant. It computes one raw relation
logit for each retained \texttt{KEY}-token and \texttt{VALUE}-token candidate
pair. V2 is the first span-level variant. It groups contiguous tokens with the
same non-\texttt{O} label before computing relation logits. V3 is a
linker-only diagnostic that reuses the V2 span-level architecture. It freezes
the encoder and entity classifier, forms spans from their predictions, and
updates only the linker. It has no separate headline benchmark result. The
final LayoutLMv3-based span-level discriminative model (V4) retains the
span-level architecture and jointly fine-tunes the encoder, classifier, and
linker.

SPADE~\cite{spade} provides context through its directed token-relation graph,
while BROS~\cite{bros} combines relative spatial representations and
area-masked pretraining with a SPADE-based relation decoder
(Section~\ref{subsec:entity_linking}). The contiguous span grouping used in
this thesis is a separate task-specific implementation decision.

Figure~\ref{fig:architecture-overview} shows the shared encoder, the
thesis-specific task modules, and the candidate paths of the development
variants. The encoder receives text tokens and two-dimensional bounding boxes
and can also accept document-image patches. The entity classifier assigns
token labels. The linker uses the encoder representations together with token
or span candidates derived from those labels. The two task heads are therefore
connected through candidate construction and are not fully independent
parallel outputs.

\begin{figure}[ht]
  \centering
  \begin{adjustbox}{max width=\textwidth}
  \begin{tikzpicture}[
    block/.style={rectangle, draw, rounded corners, minimum height=1cm, minimum width=2.8cm, align=center, font=\small},
    io/.style={rectangle, draw, minimum height=0.8cm, minimum width=2.2cm, align=center, font=\small, fill=gray!10},
    arrow/.style={-{Stealth[length=3mm]}, thick},
    node distance=0.8cm and 1.2cm
  ]
    % Input
    \node[io] (input) {Extracted Words +\\Bounding Boxes\\{\scriptsize Page image optional; not used in V4}};

    % Encoder
    \node[block, below=of input, fill=blue!10] (encoder)
      {LayoutLMv3-base Encoder\\{\scriptsize pretrained prior work; fine-tuned here}};

    % Two heads
    \node[block, below left=0.2cm and 1.5cm of encoder, fill=orange!15] (entity)
      {Entity Classifier\\\texttt{KEY} / \texttt{VALUE} / \texttt{O}\\
       {\scriptsize thesis implementation}};
    \node[block, below right=0.2cm and 1.5cm of encoder, fill=red!12] (linker)
      {Projected Scaled\\Dot-Product Linker\\with Spatial Features\\
       {\scriptsize thesis implementation}};

    % Span grouper (V2, V3, and V4)
    \node[block, below=0.8cm of entity, fill=green!12] (spans)
      {Span Grouper\\(V2/V3/V4)\\{\scriptsize thesis implementation}};

    % Output
    \node[io, below=1.35cm of encoder] (output)
      {Regular KVP Pairs\\(\texttt{KEY} $\rightarrow$ \texttt{VALUE})};

    % Arrows
    \draw[arrow] (input) -- (encoder);
    \draw[arrow] (encoder) -| (entity);
    \draw[arrow] (encoder) -| (linker);
    \draw[arrow] (entity.east) -- node[above, font=\scriptsize]
      {V1: token candidates} (linker.west |- entity.east);
    \draw[arrow] (entity) -- (spans);
    \draw[arrow] (spans) -| node[pos=0.40, above=3pt, fill=white,
      inner sep=2.5pt, font=\scriptsize\itshape, align=center]
      {V2/V4: span candidates\\V3: predicted spans} (linker);
    \draw[arrow] (linker.west |- output.east) -- (output.east);
  \end{tikzpicture}
  \end{adjustbox}
  \caption{Architecture overview and variant paths. The LayoutLMv3-base encoder
    is a pretrained component from prior work. The entity classifier, projected
    scaled-dot-product linker with spatial features, and span grouper are
    task-specific implementations developed for this thesis; the general
    scaled-dot-product principle is prior work. V1 routes token candidates
    directly to the linker. V2 and V4 use span candidates. V3 reuses the span
    path with predicted spans while keeping the encoder and classifier frozen
    and updating only the linker. During V4 training, ground-truth labels supply
    the spans; during inference, predicted labels supply them.}
  \label{fig:architecture-overview}
\end{figure}

% ---------------------------------------------------------------------------
\section{LayoutLMv3 Encoder}
\label{sec:layoutlmv3-encoder}

LayoutLMv3-base is a pretrained encoder from prior work~\cite{layoutlmv3}.
This thesis loads and fine-tunes its pretrained checkpoint. The task-specific
contributions begin with the entity classifier and relation linker described in
Sections~\ref{sec:entity-classifier} and~\ref{sec:projected-linker}.

\subsection{Input Representation}
\label{sec:input-representation}

LayoutLMv3~\cite{layoutlmv3} accepts three parallel input streams that
are fused inside the transformer backbone.

\paragraph{Text tokens.}
The extracted word sequence is tokenised with the tokenizer supplied with the
pretrained LayoutLMv3-base checkpoint~\cite{layoutlmv3}. Its vocabulary
contains 50{,}265 entries. A \texttt{[CLS]} token is prepended
and a \texttt{[SEP]} token appended. Each subword inherits the bounding box of
its parent extracted word. Multi-word entity and span boxes are formed later as the
union of their constituent word boxes.

\paragraph{Layout embeddings.}
Each token is associated with a bounding box on the normalised
$1000 \times 1000$ LayoutLMv3 coordinate grid. For text token $i$, the spatial
input is
\[
  \mathbf{b}_i =
  (x_{\min,i},\,y_{\min,i},\,x_{\max,i},\,y_{\max,i},\,w_i,\,h_i),
\]
where $\mathbf{b}_i$ is the spatial tuple for token $i$,
$(x_{\min,i},y_{\min,i})$ and $(x_{\max,i},y_{\max,i})$ are the normalised
upper-left and lower-right box coordinates,
$w_i=x_{\max,i}-x_{\min,i}$, and $h_i=y_{\max,i}-y_{\min,i}$. The x/y coordinate embeddings and width/height embeddings are combined into
the LayoutLMv3 spatial representation. The coordinate-derived components are
concatenated to form this spatial embedding, which is then combined with the
textual and sequence-position representation before the first transformer
layer.

\paragraph{Optional image patch embeddings.}
The document image is resized to $224 \times 224$ pixels and split into a
$14 \times 14$ grid of 196 non-overlapping patches. Each patch is
$16 \times 16$ pixels and is projected to the model hidden dimension
$d = 768$. Learned patch-position embeddings preserve the patch order.
The resulting 196 visual tokens are concatenated with the text tokens when
document-image input is enabled.

Although LayoutLMv3 supports document-image inputs, the reported
discriminative-model runs
did not enable page-image loading. Instead, the data pipeline passed a generated
solid-white placeholder through the LayoutLMv3 processor. This produces a
constant tensor with value $1.0$ at every pixel position, so every page receives
the same visual input. The model therefore uses text and layout information but
no document-specific visual content. This thesis did not repeat LayoutLMv3
pretraining. V4 fine-tuned the pretrained checkpoint with this constant blank
visual input and the task-specific entity and relation losses.

\subsection{Self-Attention Mechanism}
\label{sec:self-attention}

LayoutLMv3 uses a standard multi-head self-attention with $n_h = 12$ heads and
hidden dimension $d = 768$. In addition to the absolute input embeddings, the
pre-softmax attention logits include semantic one-dimensional relative-position
biases and spatial two-dimensional relative-position biases~\cite{layoutlmv3}.
The corresponding Transformers implementation exposes both bias mechanisms
and the visual patch
configuration.\footnote{\url{https://github.com/huggingface/transformers/tree/main/src/transformers/models/layoutlmv3}}
These biases let the attention mechanism represent token order and relative
document geometry.
When image patches are supplied, cross-modal attention between text tokens and
image patch tokens is unrestricted, enabling the architecture to ground textual
entities in specific image regions.

\subsection{Pre-training Objectives}
\label{sec:pretraining-objectives}

The LayoutLMv3 checkpoint used in this thesis,
\texttt{microsoft/layoutlmv3-base}, was pretrained on
IIT-CDIP~\cite{iitcdip,layoutlmv3} with three objectives:

\begin{enumerate}
  \item \textbf{\Gls{mlm}:} 30\% of text tokens are
  selected with span masking, and the model predicts the masked tokens.
  \item \textbf{\Gls{mim}:} Approximately 40\% of image
  tokens are masked, and the model reconstructs the discrete token index of
  each patch with a pretrained image tokenizer.
  \item \textbf{\Gls{wpa}:} For each unmasked word, the model
  predicts whether its associated image patches are masked. This objective
  supports alignment between text and image representations.
\end{enumerate}

These objectives define the pretrained textual, spatial, and visual
representations used to initialise the LayoutLMv3 encoder before task-specific
fine-tuning.

% ---------------------------------------------------------------------------
\section{Task Head: Entity Classifier}
\label{sec:entity-classifier}

Rather than using a sequence-to-sequence decoder, the entity head performs
token classification. It maps each text-token representation to
\texttt{KEY}, \texttt{VALUE}, or \texttt{O}. The label \texttt{O} denotes a
token that belongs to neither entity class.

\paragraph{Architecture.}
Let
\[
  \mathbf{H}
  =[\mathbf{h}_1,\ldots,\mathbf{h}_N]^\top
  \in\mathbb{R}^{N\times d}
\]
be the encoder output for $N$ text tokens, where
$\mathbf{h}_i\in\mathbb{R}^{d}$ is the representation of token $i$ and
$d=768$. The classifier first applies dropout:
\[
  \widetilde{\mathbf{h}}_i
  =\operatorname{Dropout}(\mathbf{h}_i;p),
  \qquad p=0.1.
\]
The thesis-specific linear layer then produces three raw class logits:
\[
  \mathbf{z}_i
  =\mathbf{W}_{\mathrm{ent}}\widetilde{\mathbf{h}}_i
  +\mathbf{b}_{\mathrm{ent}},
\]
where $\mathbf{z}_i\in\mathbb{R}^{3}$ is the logit vector for token $i$,
$\mathbf{W}_{\mathrm{ent}}\in\mathbb{R}^{3\times d}$ is the learned
weight matrix, and $\mathbf{b}_{\mathrm{ent}}\in\mathbb{R}^{3}$ is the
learned bias vector. Softmax gives the class-probability vector
\[
  \boldsymbol{\pi}_i=\operatorname{softmax}(\mathbf{z}_i).
\]
The predicted class is
$\widehat{y}_i=\operatorname*{arg\,max}_{c}\pi_{i,c}$ for
$c\in\{\texttt{KEY},\texttt{VALUE},\texttt{O}\}$.
Training applies cross-entropy directly to the raw logits.

\paragraph{Why token classification instead of generation?}
The discriminative and generative systems impose different output constraints.
The reconstructed Mistral-7B reference system can generate text that is not
present in the extracted input. Token classification instead assigns labels
only to observed input tokens. It therefore cannot generate new text, although
it can assign an incorrect class to an observed token.

% ---------------------------------------------------------------------------
\needspace{8\baselineskip}
\section{Projected Scaled-Dot-Product Linker (V1: Token-Level)}
\label{sec:projected-linker}

\subsection{Relation-Logit Function}
\label{sec:kv-attention}

After entity labels are available, the linker constructs candidate relations
between retained \texttt{KEY} and \texttt{VALUE} tokens or spans. It combines
one semantic compatibility scalar with an encoded geometry vector to produce
one raw relation logit for each candidate pair.

\paragraph{Semantic compatibility.}
For key candidate $i$ and value candidate $j$, let
$\mathbf{h}_i,\mathbf{h}_j\in\mathbb{R}^{d}$ be their representations.
These are token hidden states in V1 and mean-pooled span representations in the
span-level variants. The learned projections are
\[
  \mathbf{q}_i
  =\mathbf{W}_k\widetilde{\mathbf{h}}_i+\mathbf{b}_k,
  \qquad
  \mathbf{r}_j
  =\mathbf{W}_v\widetilde{\mathbf{h}}_j+\mathbf{b}_v,
\]
where $\widetilde{\mathbf{h}}_i$ and $\widetilde{\mathbf{h}}_j$ are the
candidate representations after dropout,
$\mathbf{W}_k,\mathbf{W}_v\in\mathbb{R}^{d\times d}$ are learned
projection matrices, $\mathbf{b}_k,\mathbf{b}_v\in\mathbb{R}^{d}$ are
learned bias vectors, and $d=768$. The projected scaled-dot-product
compatibility is
\[
  u_{i,j}
  =\frac{\mathbf{q}_i^\top\mathbf{r}_j}{\sqrt{d}}.
\]
The general scaled-dot-product principle follows
Vaswani et al.~\cite{vaswani2017attention}.

\paragraph{Design evolution.}
An earlier implementation used a full biaffine relation function with a
learned bilinear interaction matrix, following Dozat and
Manning~\cite{dozat2017biaffine}. It produced a not-a-number
(\texttt{NaN}) link loss at approximately step~1400. The later projected
scaled-dot-product implementation did not show this failure in the subsequent
reported runs. The active V4 component is a span-level projected
scaled-dot-product linker with spatial features. Its Python class keeps its
historical name only for checkpoint compatibility.

\subsection{Spatial Encoding}
\label{sec:spatial-encoding}

For candidate boxes $i$ and $j$, the implementation constructs the
eight-dimensional geometry vector
\[
  \mathbf{x}_{i,j}
  =(\Delta x_{i,j},\Delta y_{i,j},r_{i,j},\phi_{i,j},
    o_{y,i,j},o_{x,i,j},\rho_{A,i,j},\rho_{R,i,j}).
\]
Here, $\Delta x_{i,j}$ and $\Delta y_{i,j}$ are the value-box centre
coordinates minus the key-box centre coordinates,
$r_{i,j}$ is their Euclidean distance, and
$\phi_{i,j}=\operatorname{atan2}(\Delta y_{i,j},\Delta x_{i,j})$ is their
angle. The terms $o_{y,i,j}$ and $o_{x,i,j}$ are the overlaps of the box
intervals on the y and x axes, normalised by the smaller box height and width,
respectively. The term $\rho_{A,i,j}$ is the value-to-key area ratio, and
$\rho_{R,i,j}$ is the value-to-key aspect-ratio ratio.

A two-layer \gls{mlp} with rectified linear unit (ReLU) activations maps
$\mathbf{x}_{i,j}$ to a spatial representation
$\mathbf{g}_{i,j}\in\mathbb{R}^{32}$. A final MLP produces
\[
  \ell_{i,j}
  =f_{\mathrm{out}}([u_{i,j};\mathbf{g}_{i,j}]),
\]
where $[\cdot\,;\cdot]$ denotes concatenation,
$f_{\mathrm{out}}$ is the learned final MLP, and
$\ell_{i,j}\in\mathbb{R}$ is the raw relation logit. The sigmoid function
converts this logit to a relation probability:
\[
  p_{i,j}
  =\sigma(\ell_{i,j})
  =\frac{1}{1+\exp(-\ell_{i,j})}.
\]
Training uses the raw relation logit in the
binary-cross-entropy-with-logits loss. Inference applies a threshold to the
relation probability. Neither quantity is a benchmark evaluation metric.
Chapter~\ref{ch:implementation} records the exact numerical safeguards used
by the implementation.

\subsection{Class Imbalance Problem}
\label{sec:v1-imbalance}

Let $N_K$ and $N_V$ be the numbers of retained \texttt{KEY} and
\texttt{VALUE} token candidates on one page. V1 computes $N_KN_V$ relation
logits, arranged as an $N_K\times N_V$ matrix. The matrix unit is a token-pair
candidate, not an evaluated benchmark pair. For illustration, a page with 100
key-token candidates and 100 value-token candidates produces 10{,}000
candidate pairs. If 20 pairs are positive, the negative-to-positive ratio is
approximately \textbf{500:1}. Despite applying a positive-class weight to the
\gls{bce} loss, this extreme imbalance was associated with poor relation
learning in the validation diagnostics.

This observation motivates span-level linking. An illustrative case with 15
key spans and 15 value spans produces 225 candidate pairs. With 15 positive
relations, it has 210 negative candidates and a negative-to-positive ratio of
approximately $14{:}1$.

% ---------------------------------------------------------------------------
\section{Span-Level Linker (V2/V3/V4)}
\label{sec:v2-span-linking}

V2 is the first span-level development variant. It reduces the granularity
mismatch of V1 by linking contiguous entity spans instead of individual
tokens. This changes the candidate unit from a token pair to a
key-span--value-span pair. The contiguous-span construction is specific to
this thesis implementation. SPADE~\cite{spade} and BROS~\cite{bros} provide
broader context for learned token relations and spatially informed
representations, respectively.

\subsection{Span Grouping}
\label{sec:span-grouping}

Contiguous tokens with the same non-\texttt{O} entity label are merged into
spans. For example, if tokens 5--7 are labelled
\texttt{KEY} and tokens 8--12 are labelled \texttt{VALUE}, they form one key
span and one value span, respectively.

The first V2 training path formed spans from predicted entity labels while
ground-truth annotations supplied the relation labels. Entity-prediction
errors could therefore misalign the candidate spans and relation labels. A
later V2 training configuration formed spans from ground-truth entity labels.
This configuration is referred to here as teacher-forced span construction.

V3 was a linker-only diagnostic that loaded a teacher-forced V2 checkpoint.
It froze the LayoutLMv3 encoder and entity classifier and fine-tuned only the
linker. Freezing these components keeps the upstream encoder and classifier
parameters fixed, so parameter updates in V3 are confined to relation linking
while candidates are formed from predicted spans. During V3 training, spans
came from the frozen classifier predictions, while ground-truth relation labels
supplied the linker targets. V3 was not a separate final architecture and has
no headline benchmark result.

V4 uses ground-truth entity labels for span construction during training and
predicted labels during inference. Unlike V3, V4 jointly optimises the encoder,
entity classifier, and linker. It therefore retains a train--inference
difference in the source of its spans.

Figure~\ref{fig:span-grouping} illustrates the span grouping process.

\begin{figure}[ht]
  \centering
  \begin{adjustbox}{max width=\textwidth}
  \begin{tikzpicture}[
    tok/.style={rectangle, draw, minimum height=0.7cm, minimum width=1.2cm,
                align=center, font=\scriptsize},
    span/.style={rectangle, draw, rounded corners, minimum height=0.8cm,
                 align=center, font=\scriptsize\bfseries},
    arrow/.style={-{Stealth[length=2mm]}, thick, gray},
    node distance=0.05cm
  ]
    % Token row
    \node[tok, fill=gray!15] (t1) {Invoice};
    \node[tok, fill=red!20, right=of t1] (t2) {Date};
    \node[tok, fill=red!20, right=of t2] (t3) {:};
    \node[tok, fill=blue!20, right=of t3] (t4) {01};
    \node[tok, fill=blue!20, right=of t4] (t5) {/15/};
    \node[tok, fill=blue!20, right=of t5] (t6) {2024};
    \node[tok, fill=red!20, right=of t6] (t7) {Total};
    \node[tok, fill=red!20, right=of t7] (t8) {Amt};
    \node[tok, fill=blue!20, right=of t8] (t9) {\$1500};

    % Labels
    \node[font=\tiny, above=0.15cm of t1] {\texttt{O}};
    \node[font=\tiny, above=0.15cm of t2] {\texttt{KEY}};
    \node[font=\tiny, above=0.15cm of t3] {\texttt{KEY}};
    \node[font=\tiny, above=0.15cm of t4] {\texttt{VALUE}};
    \node[font=\tiny, above=0.15cm of t5] {\texttt{VALUE}};
    \node[font=\tiny, above=0.15cm of t6] {\texttt{VALUE}};
    \node[font=\tiny, above=0.15cm of t7] {\texttt{KEY}};
    \node[font=\tiny, above=0.15cm of t8] {\texttt{KEY}};
    \node[font=\tiny, above=0.15cm of t9] {\texttt{VALUE}};

    % Span row
    \node[span, fill=red!30, below=1.2cm of $(t2)!0.5!(t3)$] (s1) {Key 1:\\``Date :''};
    \node[span, fill=blue!30, below=1.2cm of t5] (s2) {Value 1:\\``01/15/2024''};
    \node[span, fill=red!30, below=1.2cm of $(t7)!0.5!(t8)$] (s3) {Key 2:\\``Total Amt''};
    \node[span, fill=blue!30, below=1.2cm of t9] (s4) {Value 2:\\``\$1500''};

    % Arrows
    \draw[arrow] (t2.south) -- (s1.north);
    \draw[arrow] (t3.south) -- (s1.north);
    \draw[arrow] (t4.south) -- (s2.north);
    \draw[arrow] (t5.south) -- (s2.north);
    \draw[arrow] (t6.south) -- (s2.north);
    \draw[arrow] (t7.south) -- (s3.north);
    \draw[arrow] (t8.south) -- (s3.north);
    \draw[arrow] (t9.south) -- (s4.north);
  \end{tikzpicture}
  \end{adjustbox}
  \caption{Span grouping: contiguous tokens with the same entity label are
    merged into spans. The span representation is the mean pool of constituent
    token hidden states; the span bounding box is the union of constituent
    token boxes.}
  \label{fig:span-grouping}
\end{figure}

\subsection{Span Representation}
\label{sec:span-representation}

For span $k$, let $T_k$ be the set of indices of its constituent tokens. Its
mean-pooled representation is
\[
  \mathbf{s}_k
  =\frac{1}{|T_k|}
   \sum_{i\in T_k}\mathbf{h}_i,
\]
where $\mathbf{s}_k\in\mathbb{R}^{d}$ is the span representation,
$|T_k|$ is the number of tokens in the span, and
$\mathbf{h}_i\in\mathbb{R}^{d}$ is the encoder representation of token $i$.
The span bounding box is the tight union box over all constituent token boxes.
The same spatial
augmentation (8 geometric features, Section~\ref{sec:spatial-encoding}) is
applied at the span level.

\subsection{Span-Level Relation Logits}
\label{sec:span-scoring}

The span-level projected scaled-dot-product linker operates identically to V1
(Section~\ref{sec:kv-attention}) but over span representations instead of
token representations. For $K$ retained key spans and $V$ retained value
spans, the linker produces a relation-logit matrix
$\mathbf{L}\in\mathbb{R}^{K\times V}$, where entry
$L_{i,j}=\ell_{i,j}$ is the raw relation logit for key span $i$ and value span
$j$. For illustration, 15 key spans and 15
value spans produce 225 candidate pairs. With 15 positives and 210 negatives,
this gives a negative-to-positive ratio of approximately \textbf{$14{:}1$},
compared with approximately $500{:}1$ in the illustrative V1 example.

\subsection{Relation Supervision}
\label{sec:v2-training}

The span-level variants use the loss definition in
Section~\ref{sec:loss-function}. Their candidate unit changes from a token
pair in V1 to a span pair in V2, V3, and V4. The binary relation target remains
one for a linked candidate and zero otherwise.

For V4 relation supervision, each recovered ground-truth KVP contributes one
representative word-level \texttt{KEY}$\rightarrow$\texttt{VALUE} relation.
V4 does not use the earlier key-word $\times$ value-word cross-product.
Tokenisation can expand the representative relation into multiple positive
subword-level cells before the labels are collapsed to the span level.
Chapter~\ref{ch:implementation} describes the implementation details.

% ---------------------------------------------------------------------------
\section{Training Procedure}
\label{sec:training-procedure}

\subsection{Loss Function}
\label{sec:loss-function}

For the jointly trained variants, the total objective is
\begin{equation}
  \label{eq:joint-loss}
  \mathcal{L}
  =\mathcal{L}_{\mathrm{entity}}
  +\lambda\mathcal{L}_{\mathrm{link}},
\end{equation}
where $\mathcal{L}$ is the total loss,
$\mathcal{L}_{\mathrm{entity}}$ is the entity-classification loss,
$\mathcal{L}_{\mathrm{link}}$ is the relation loss, and $\lambda$ controls
the relation-loss contribution.

Let $\mathcal{A}$ be the set of active, non-padding text-token positions,
$y_i$ the ground-truth class of token $i$, and $\pi_{i,y_i}$ the predicted
probability of that class. The unweighted entity loss is
\[
  \mathcal{L}_{\mathrm{entity}}
  =-\frac{1}{|\mathcal{A}|}
   \sum_{i\in\mathcal{A}}\log\pi_{i,y_i},
  \qquad
  y_i\in\{\texttt{KEY},\texttt{VALUE},\texttt{O}\}.
\]
No entity-class weight is applied.

For page $b$, let $\mathcal{C}_b$ be its set of valid binary relation
candidates and let $y_{b,i,j}\in\{0,1\}$ be the ground-truth relation label
for candidate pair $(i,j)$. The positive-class weight for page $b$ is
\[
  w_b^{+}
  =\min\!\left(
    \frac{\max(n_{\mathrm{neg},b},1)}
         {\max(n_{\mathrm{pos},b},1)},
    50
  \right),
\]
where $n_{\mathrm{pos},b}$ and $n_{\mathrm{neg},b}$ are the numbers of
positive and negative valid relation labels. The page-level relation loss is
\[
  \mathcal{L}_{\mathrm{link}}^{(b)}
  =-\frac{1}{|\mathcal{C}_b|}
   \sum_{(i,j)\in\mathcal{C}_b}
   \left[
     w_b^{+}y_{b,i,j}\log\sigma(\ell_{b,i,j})
     +(1-y_{b,i,j})
      \log\!\left(1-\sigma(\ell_{b,i,j})\right)
   \right],
\]
where $\ell_{b,i,j}$ is the raw relation logit. Let
$\mathcal{B}_{\mathrm{link}}$ be the set of batch pages with at least one
valid relation candidate. The batch relation loss is
\[
  \mathcal{L}_{\mathrm{link}}
  =\frac{1}{|\mathcal{B}_{\mathrm{link}}|}
   \sum_{b\in\mathcal{B}_{\mathrm{link}}}
   \mathcal{L}_{\mathrm{link}}^{(b)}.
\]
The implementation uses the numerically stable
binary-cross-entropy-with-logits operation rather than calculating the
logarithms explicitly. Earlier development runs tested
$\lambda\in\{0.5,1.0,2.0\}$; V4 uses $\lambda=5.0$. V3 optimises only the
weighted relation loss because its encoder and entity classifier are frozen.

\paragraph{Training stability.}
The reported joint-training implementation uses projected scaled-dot-product
relation logits (Section~\ref{sec:kv-attention}) and applies a valid-target
mask before the \gls{bce}-with-logits operation. The exact numerical
safeguards and low-level loss procedures are given in
Section~\ref{sec:linker-impl}. The reported 32-bit floating-point (fp32) runs
did not show the same \texttt{NaN} failure. The validation-based diagnostic
analysis is reported in Section~\ref{sec:exp567}.

\subsection{Optimisation}
\label{sec:optimisation}

We use AdamW~\cite{loshchilov2019decoupled} with weight decay $0.01$.
The final V4 uses batch size 1, gradient accumulation over 8 batches, a unified
learning rate of $2\times10^{-5}$, and a maximum of 30 epochs. Early stopping
is permitted from epoch 10 and uses patience 5. The scheduler is defined over
actual optimiser updates. With 4{,}851 training batches, one epoch contains
$\lceil4851/8\rceil=607$ updates. The linear schedule therefore has at most
18{,}210 optimiser steps, with the first 500 used for warm-up; early stopping
can reduce the completed total. The V4 decoder uses a fixed
relation-probability threshold of $0.5$ when producing key--value pairs. The
released KVP10k evaluator uses \gls{ned} $<0.2$ and \gls{iou} $\geq0.3$ for
text and location matching. For V4, checkpoint selection and early stopping
use only the validation Regular text+location macro pair F1; test metrics are
not used for these decisions. The validation metric selects checkpoint 15.
Training completes 20 epochs and 12{,}140 optimiser steps before early stopping.

An earlier bfloat16 development attempt produced a \texttt{NaN} loss. All
reported discriminative-model runs therefore use fp32 precision.

\subsection{Development Initialisation and Joint Fine-Tuning}
\label{sec:two-phase}

The final V4 training path contains an entity-classification initialisation phase
and a subsequent joint entity-and-linker phase. These phases do not form one
continuous run.

\begin{enumerate}
  \item \textbf{Entity-classification initialisation:}
  The encoder and entity-classification head are trained with $\lambda=0$,
  which disables the relation objective. Validation class-aware
  \texttt{KEY}/\texttt{VALUE} micro F1 selects epoch~10. Test data is not
  accessed during this phase.

  \item \textbf{Joint entity-and-linker fine-tuning:}
  V4 loads the selected entity encoder and classifier and creates a fresh,
  randomly initialised relation linker. It then jointly optimises all three
  components and does not reuse the linker from an earlier joint checkpoint.
  The final phase uses span-level candidate construction, representative
  relation supervision, class-aware validation metrics, and official validation
  Regular text+location macro pair-F1 checkpoint selection. This metric selects
  checkpoint~15. Chapter~\ref{ch:implementation} records the exact checkpoint
  provenance.
\end{enumerate}

% ---------------------------------------------------------------------------
\section{Inference}
\label{sec:inference}

\subsection{Entity Extraction}
\label{sec:greedy-decoding}

At inference time, softmax converts each per-token class-logit vector
$\mathbf{z}_i$ to the probability vector $\boldsymbol{\pi}_i$. The class with
the largest probability is selected from
$\{\texttt{KEY},\texttt{VALUE},\texttt{O}\}$. Consecutive tokens with the
same non-\texttt{O} label are then merged into spans. The bounding box of a
span is the tight union box over all constituent token boxes.

For every retained candidate pair, the linker produces a raw relation logit
$\ell_{i,j}$. Sigmoid converts each logit to a relation probability
$p_{i,j}=\sigma(\ell_{i,j})$. For each key candidate $i$, the decoder selects
\[
  j_i^\ast
  =\operatorname*{arg\,max}_{j}p_{i,j}
  =\operatorname*{arg\,max}_{j}\ell_{i,j}.
\]
The equality holds because sigmoid is monotonic. The decoder emits the relation
$i\rightarrow j_i^\ast$ if $p_{i,j_i^\ast}\geq\tau$, where the fixed
threshold is $\tau=0.5$. Otherwise, the key candidate remains unmatched. Thus,
each key candidate produces at most one predicted relation.

The thesis implementation serialises these thresholded predictions for the
released KVP10k evaluator. This integration is thesis work. The evaluator
itself is a prior resource and computes benchmark precision, recall, and F1
from the emitted predictions.

\subsection{Linking (V1)}
\label{sec:v1-linking-inference}

In V1, $i$ and $j$ refer to retained \texttt{KEY} and \texttt{VALUE} tokens.
The decoder applies the common probability rule above to the
$N_K\times N_V$ relation-logit matrix.

\subsection{Linking (V2/V3/V4)}
\label{sec:v2-linking-inference}

In V2, V3, and V4, $i$ and $j$ refer to retained key and value spans. At
inference time, the spans are formed from predicted entity labels and are
represented by mean-pooled encoder states. The decoder applies the common
probability rule above to the $K\times V$ relation-logit matrix.

\subsection{Post-processing}
\label{sec:post-processing}

The final V4 evaluation uses the learned direct decoder. It emits Regular
linked key--value span pairs whose selected relation probability is at least
$0.5$. Unmatched spans are not serialised, so the final output contains no
Unkeyed or Unvalued predictions. The separate recovery method from earlier
development is not part of the checkpoint-15 result.
No beam search or constrained decoding is applied, and bounding boxes are
inherited directly from the extracted input.
